% ── Section II: Related Work ─────────────────────────────────────────
% Conference budget: ~0.75 page. Three threads + positioning paragraph.
% TODO before submission: add 1–2 recent (2024–2026) accident-detection
% and adaptive-perception citations per thread; see references.bib.
% ─────────────────────────────────────────────────────────────────────

\section{Related Work}
\label{sec:related}

\subsection{Accident Detection in Traffic Surveillance}
\label{subsec:rw_accident}

Vision-based accident detection is commonly formulated as clip- or
frame-level classification over surveillance footage.
The CADP dataset~\cite{Shah2018_cadp} established CCTV-based accident
analysis; D2-City~\cite{Che2019_d2city} provides large-scale dashcam
footage of diverse (including Chinese urban) traffic.
Deployed commercial systems such as Hikvision AID~\cite{Hikvision_AID}
detect incidents on highways with car-centric assumptions and offer no
severity grading.
Across this line of work, classifiers run with fixed operating points
tuned offline on pre-segmented clips; none adapts its thresholds to the
traffic state it currently observes, which is precisely the regime in
which false alerts accumulate in dense mixed traffic.

\subsection{Multi-Object Tracking in Urban Scenes}
\label{subsec:rw_mot}

Tracking-by-detection dominates modern MOT.
SORT~\cite{Bewley2016_sort} paired a linear Kalman filter with
IoU-based Hungarian association; ByteTrack~\cite{Zhang2022_bytetrack}
improved recall by also associating low-score detections;
StrongSORT~\cite{Du2023_strongsort} strengthened appearance cues; and
recent occlusion-aware trackers exploit scene
structure~\cite{Hashempoor2025_fasttracker}.
Benchmarks such as MOT20~\cite{Dendorfer2020_mot20} focus on
pedestrians and cars, whose image-plane motion is well served by a
constant-velocity assumption.
All of these trackers apply one linear motion model to every object
class.
Nonlinear models with an explicit turn-rate state, notably CTRV,
were shown to outperform constant-velocity models for maneuvering
vehicles in the sensor-fusion literature~\cite{Schubert2008_ctrv}, yet
class-conditioned motion modeling --- and motorcycle weaving in
particular --- remains unaddressed in tracking-by-detection MOT.

\subsection{Context-Aware and Adaptive Perception}
\label{subsec:rw_adaptive}

Adaptive perception work mostly targets environmental domain shift:
Image-Adaptive YOLO~\cite{Liu2022_iayolo} learns image-enhancement
parameters for adverse weather, and related methods adjust
preprocessing or confidence calibration for night or rain.
These approaches adapt to \emph{appearance} conditions but leave the
tracker's hyperparameters and motion assumptions fixed, and they
optimize a single downstream objective.
Serving two tasks with conflicting operating points --- precision-first
alerting and recall-first control observations, as arise when
perception feeds both a safety channel and a multi-agent controller
trained with MAPPO~\cite{Yu2022_mappo} --- has not been formulated at
the perception level.

\medskip
\noindent\textbf{Positioning.}
To the best of our knowledge, no prior system simultaneously provides
(i)~perception-level parameter adaptation driven by the observed
traffic state, (ii)~class-conditioned motion modeling for
motorcycle-dominant mixed traffic, and (iii)~a single track-quality
formulation serving two downstream tasks with different
precision--recall requirements.
The proposed pipeline addresses all three.
